\subsection{DISTRIBUTIVITY OF ONE INTERNAL LAW WITH RESPECT TO ANOTHER}

DEFINITION 7. Let \(\top\) and \(\bot\) be two internal laws on a set E. The law \(\bot\) is said to be distributive with respect to the law \(\top\) if

\begin{enumerate}
\def\labelenumi{(\arabic{enumi})}
\setcounter{enumi}{9}
\tightlist
\item
\end{enumerate}

\[
x \perp (y \top z) = (x \perp y) \top (x \perp z)\tag{11}
\]

\[
(x \top y) \perp z = (x \perp z) \top (y \perp z)
\]

for all \(x, y, z\) in \(\mathbf{E}\) .

Note that (10) and (11) are equivalent if the law \(\perp\) is commutative. In general, one of the laws is written additively and the other multiplicatively; if multiplication is distributive with respect to addition, then:

\begin{enumerate}
\def\labelenumi{(\arabic{enumi})}
\setcounter{enumi}{11}
\tightlist
\item
\end{enumerate}

\[
x. (y + z) = x. y + x. z\tag{13}
\]

\[
(x + y). z = x. z + y. z
\]

Examples. (1) In the set \(\mathfrak{P}(\mathbf{E})\) of subsets of a set \(\mathbf{E}\) , each of the internal laws \(\cap\) and \(\cup\) is distributive with respect to itself and the other. This follows from formulae of the form

\[
\begin{array}{l} \mathbf {A} \cap (\mathbf {B} \cup \mathbf {C}) = (\mathbf {A} \cap \mathbf {B}) \cup (\mathbf {A} \cap \mathbf {C}) \\ \mathbf {A} \cup (\mathbf {B} \cap \mathbf {C}) = (\mathbf {A} \cup \mathbf {B}) \cap (\mathbf {A} \cup \mathbf {C}). \end{array}
\]

\begin{enumerate}
\def\labelenumi{(\arabic{enumi})}
\setcounter{enumi}{1}
\item
  In Z (and more generally, in any totally ordered set) each of laws sup and inf is distributive with respect to the other and with respect to itself.
\item
  In \(\mathbf{Z}\) (*and more generally in any ring*) multiplication is distributive with respect to addition.
\item
  In N addition and multiplication are distributive with respect to the laws sup and inf.
\end{enumerate}

